% conclusion.tex
% Three short paragraphs: contribution recap, validation summary,
% future work.

We presented a two-stage construction for synthetic equity path
generation that closes two gaps in the literature: a hybrid hidden
Markov framework that reproduces single-asset stylized facts via a
Laplace quantile partition and a Poisson jump-duration override, and
a closed-form variance-corrected single-index composer that places
those validated marginals in a multi-asset factor structure without
inflating the marginal variance. The construction is
generator-agnostic on the composition side, in the sense that any
per-asset full-return generator can plug into the variance correction
and inherit the marginal it was built to deliver.

On the empirical side, the hybrid HMM marginals passed the
Kolmogorov-Smirnov test against the real SPY univariate distribution
at $97.6\%$ in-sample and $94.4\%$ out-of-sample, outperforming
Bootstrap, Gaussian, Laplace, GARCH, hidden semi-Markov, and Gated
Recurrent Unit baselines on combined distributional and tail metrics.
The variance-corrected single-index composer recovered the factor
loading to a median absolute error of $0.018$ and matched the leading
residual-fit alternatives on a per-ticker one-day $99\%$ Value-at-Risk
backtest, validating the composed paths against the real $2014$--$2024$
growth-rate histories. We documented one downstream caveat: an
iid-residual diversification-return artifact under daily rebalancing
that inflates absolute return levels by roughly $22$ percentage points
on a deployed $20$-asset target-weight allocator. A uniform
location-shift bias correction anchored to a documented long-run
prior addresses the artifact as a per-path monotone transform that
preserves rank order, drawdown shape, and volatility envelope while
shifting only the level.

Several extensions follow from the present work. The
variance correction generalizes to multi-factor and lagged-factor
single-index models without changing the algebra; we leave
the empirical evaluation of these extensions to follow-up work.
Per-ticker tuning of the jump parameters $(\epsilon, \lambda)$ would
push the hybrid composer's marginal-fidelity pass rates upward at the
cost of a calibration step per asset. The diversification-return
artifact admits a structural fix via autoregressive or block-bootstrap
residuals composed with the single-index factor; that fix requires
restructuring the copula sampling step and is the most consequential
follow-up for high-turnover downstream workflows.
